\documentclass[11pt, a4paper]{article}

\usepackage[T1]{fontenc}
\usepackage[utf8]{inputenc}
\usepackage{tgpagella}
\usepackage[spanish, es-noshorthands]{babel}
\usepackage{amsmath, amssymb}
\usepackage{booktabs}
\usepackage{tabularx}
\usepackage[most]{tcolorbox}
\usepackage[margin=2.5cm]{geometry}
\usepackage{fancyhdr}

\pagestyle{fancy}
\fancyhf{}
\setlength{\headheight}{16pt}
\lhead{\textbf{Ing. Mecatrónica}}
\chead{Preparación Asincrónica: Lab 4}
\rhead{Circuitos Electrónicos II}
\renewcommand{\headrulewidth}{0.4pt}
\definecolor{ucbblue}{RGB}{0, 51, 102}

\begin{document}

\begin{center}
    \Large\textbf{\textcolor{ucbblue}{Preparación Anticipada: Limitaciones Dinámicas (Lab 4)}}[cite: 11]
\end{center}

\begin{tcolorbox}[colback=white, colframe=black, title=\textbf{Objetivo Experimental[cite: 11]}]
    ¿Por qué un amplificador que funciona correctamente a DC o baja frecuencia puede distorsionar una señal cuando aumentan la frecuencia o la amplitud?[cite: 11]
\end{tcolorbox}

\section*{1. Identificación de Dispositivos[cite: 11]}
\textbf{Dispositivo 1:} \rule{4cm}{0.4pt} \textbf{¿LM741 confirmado?} Sí / No[cite: 11] \\
\textbf{Dispositivo 2:} \rule{4cm}{0.4pt} \textbf{¿TL082 confirmado?} Sí / No[cite: 11] \\
\textit{No utilizar estos modelos como definitivos hasta confirmación expresa del instructor.[cite: 11]}

\section*{2. Extracción de Datasheet[cite: 11]}
\begin{table}[h]
    \centering
    \renewcommand{\arraystretch}{1.5}
    \begin{tabularx}{\textwidth}{|l|X|X|}
    \hline
    \textbf{Campo / Especificación} & \textbf{Dispositivo 1} & \textbf{Dispositivo 2} \\ \hline
    Fabricante[cite: 11] & & \\ \hline
    Número de parte exacto[cite: 11] & & \\ \hline
    Revisión del Datasheet[cite: 11] & & \\ \hline
    GBW (Ganancia Ancho de Banda)[cite: 11] & & \\ \hline
    GBW Condiciones de prueba[cite: 11] & & \\ \hline
    Slew Rate (SR)[cite: 11] & & \\ \hline
    SR Condiciones de prueba[cite: 11] & & \\ \hline
    Alimentación usada en especificación[cite: 11] & & \\ \hline
    Notas adicionales[cite: 11] & & \\ \hline
    \end{tabularx}
\end{table}

\section*{3. Cuestionario Conceptual[cite: 11]}
Responda brevemente fundamentando en la teoría de circuitos dinámicos:
\begin{enumerate}
    \item ¿Qué significa físicamente que el ancho de banda de un Op-Amp sea finito?[cite: 11]
    \item ¿Por qué la configuración de ganancia en lazo cerrado puede afectar la frecuencia máxima útil del circuito?[cite: 11]
    \item ¿Qué representa matemáticamente y físicamente la expresión $SR = \max\left|\frac{dV_o}{dt}\right|$?[cite: 11]
    \item Para una señal senoidal $v_o(t) = V_p\sin(2\pi ft)$, ¿qué ocurre con la rapidez máxima requerida (derivada) cuando aumenta la amplitud $V_p$?[cite: 11] ¿Y cuando aumenta la frecuencia $f$?[cite: 11]
    \item ¿Qué evidencia experimental permitiría diferenciar claramente una limitación dinámica de un simple error DC o saturación estática?[cite: 11]
\end{enumerate}

\end{document}
